\chapter*{Resumo}
\pagenumbering{gobble}
\noindent Grande parte da informação clínica de pacientes com doenças hepáticas está registrada em texto livre nos prontuários eletrônicos, o que torna sua extração manual lenta e sua análise sistemática difícil. Modelos de linguagem de grande porte podem automatizar essa extração, mas, combinados à recuperação tradicional por similaridade, têm dificuldade com evidências dispersas ao longo do texto. Este trabalho avalia o GraphRAG, abordagem que combina esses modelos a grafos de conhecimento, na identificação de oito doenças hepáticas em prontuários eletrônicos reais, em português brasileiro, do Hospital Universitário da Universidade Federal de Alagoas. Sobre uma amostra de 537 prontuários anotados manualmente, foram comparados os modos de busca \textit{basic}, \textit{local} e \textit{drift} em dois níveis de comunidade do grafo. O modo \textit{drift} obteve o melhor desempenho (F1 macro de 0,798), graças a uma revocação muito mais alta, ao custo de menor precisão. O ganho foi maior nas doenças cuja evidência está dispersa no texto, como cirrose e hipertensão portal, e pequeno ou nulo nas que possuem marcador textual direto; o nível de comunidade influenciou menos que o modo de busca. Conclui-se que a escolha do modo de busca deve considerar o tipo de evidência textual de cada doença.

\vspace{5mm}
\noindent\textbf{\textit{Palavras-chave}:~}GraphRAG; processamento de linguagem natural; prontuário eletrônico; hepatologia; classificação multirrótulo.
